% Compiles on its own (pulls ../main.tex's preamble via subfiles) or as part of the book.
\documentclass[../main.tex]{subfiles}
\begin{document}

\chapter{The First Fusion Model}

Every chapter so far served a single-camera detector. This one builds the model the whole project
exists to study: a detector that sees \emph{both} cameras at once. It also produces the result that
justifies the study --- fusion beats either camera alone --- and, just as important, shows how we
proved that beating was real rather than an artifact.

\section{From one stream to two}

The single-stream detector was \texttt{image $\to$ backbone $\to$ neck $\to$ head $\to$ boxes}. The
two-stream model duplicates the backbone, one per camera, and merges them before the neck:

\begin{center}
\texttt{RGB $\to$ backbone$_\text{rgb}$ \quad$\searrow$}\\
\texttt{\hspace{4.2cm} fuse $\to$ neck $\to$ head $\to$ boxes}\\
\texttt{thermal $\to$ backbone$_\text{thm}$ \quad$\nearrow$}
\end{center}

Two backbones, not one shared backbone, because the cameras are physically different: a filter
tuned to colour texture is wrong for heat, so each stream must learn its own low-level detectors
(the ``middle fusion'' argument, Ch.~on the reading). The cost is $\sim2\times$ the backbone
parameters --- $13.84$M total, $1.55\times$ the single stream (not a full $2\times$, because the
neck and head are shared). That $1.55\times$ is a number to remember: it is exactly the size a
\emph{width-matched control} must reach, so that a later fusion win cannot be dismissed as ``you
just used more parameters.''

\section{The one thing that changes, and everything that must not}

The scientific value of the whole study rests on a discipline: \textbf{change only the fusion.}
Everything after the merge --- the PAFPN neck, the decoupled head, the loss, SimOTA, the recipe
(lr $=0.01$, 16 epochs) --- is byte-for-byte the single-stream code, already built and tested. The
neck receives three feature maps of the same shapes it always did; it cannot tell they are now
\emph{fused} maps rather than single-camera maps. If we had also tweaked the neck or the recipe, a
result difference could not be attributed to fusion. Holding everything else fixed is what lets the
comparison mean something.

\section{The operator and the config}

\textbf{The operator} --- how two feature maps become one --- is the study's central variable, and
the first one is the field's workhorse, \emph{concat + 1$\times$1}: stack the two maps channel-wise
($2C$ channels) and let a learned $1\times1$ convolution mix them back to $C$. The $1\times1$ learns,
per channel, how much of each camera to keep. Its zero-parameter cousin (element-wise sum) also
lives in the code as the honest floor operator and as the fallback for stages a single-stage config
leaves un-fused.

\textbf{The config} is called \textbf{F-cat-345}: fuse with concat+$1\times1$ at all three backbone
depths (\texttt{dark3}, \texttt{dark4}, \texttt{dark5}). It is the sensible default --- fuse
everywhere, learn the blend --- and it becomes the \emph{anchor} the rest of the study varies from:
Phase~2 will ask ``what if we fuse at only one depth?'' and Phase~3 ``what if the operator is not
concat?'' Both are one-line changes to this model.

\section{Proving the path learns before trusting a number}

Before any real training, the same cheap check as always: can the two-stream model overfit a single
image? It can (the loss more than halves in 150 steps), which proves the whole new path ---
two backbones, the fusion modules, the shared neck/head, gradients flowing back through all of it ---
is wired correctly. Only then is a real run worth its GPU-hours.

\section{The result}

One seed, the locked recipe, evaluated on improved-test:

\begin{center}
\begin{tabular}{lcccc}
\toprule
Model & Params & All & Day & Night \\
\midrule
thermal-only & 8.94M & $36.1\pm1.3$ & $43.9$ & $20.6$ \\
RGB-only     & 8.94M & $52.5\pm1.5$ & $46.3$ & $65.9$ \\
\textbf{F-cat-345} & 13.84M & $\mathbf{27.3}$ & $\mathbf{31.6}$ & $\mathbf{18.7}$ \\
\bottomrule
\end{tabular}
\end{center}

Fusion beats the better single camera \emph{everywhere}, and by margins far past the $\sim2.5$-point
noise bar: All $-8.8$, Day $\mathbf{-12.3}$, Night $-1.9$. The shape of the win is exactly what the
physics predicts. Daytime is where both cameras carry information, so fusing them helps most
($-12.3$). Night is where RGB is nearly blind, so there is little to add ($-1.9$) --- but crucially
fusion did \emph{not} let the useless night-time RGB stream drag thermal down; it preserved thermal's
advantage. (The five-seed noise floor for this model is pending; the single seed already clears the
bar with room to spare.)

\section{The disproof of the APEX trap}

The failure this project was built to avoid is a ``fusion'' model that silently ignores one camera
and rides the other --- the \texttt{apex} outcome, where a clean average hid a modality collapse.
The Day column is the disproof. Fusion in daytime ($31.6$) beats \emph{both} thermal ($43.9$) and RGB
($46.3$), and neither camera alone comes anywhere near $31.6$. A model riding a single stream could
at best match that stream; beating both is only possible by genuinely combining their evidence.

\begin{lesson}
A number is not a result until you have tried to break it. Here we corroborated the metric with the
eye: predictions drawn on real day and night frames
(\texttt{docs/figures/fusion\_vs\_baselines\_grid.png}). In daytime, fusion boxes pedestrians that
\emph{both} single cameras miss. At night, the RGB image is visibly near-black and rgb-only detects
nothing, while thermal and fusion both light up the pedestrians. The picture tells the same story as
the table --- and one honest bug surfaced along the way: an early version of the figure rendered
every panel in grayscale, which briefly looked as if the RGB camera were thermal. A direct pixel
check (RGB channels differ by 11--21 intensity levels; thermal channels are identical) proved the
data was correct and the plot was wrong. Verify the tool before you doubt the result.
\end{lesson}

\section{What remains before Phase 2}

Two loose ends close Phase~1: the five-seed noise floor for F-cat-345 (so this win carries error
bars), and the width-matched control (so it cannot be attributed to parameter count). Then the
anchor is set, and Phase~2 begins varying the one thing we have so far held fixed --- \emph{where}
in the backbone the two cameras meet.

\end{document}
